\documentclass[10pt,a4paper]{amsart}
\usepackage{amsmath,amssymb,amsthm,graphicx}
% Retained preamble for the Stacks Project chapter divisors.tex (Divisors).
% The chapter file pinned-original-source.tex carries only \input{preamble} on its first line, so no
% preamble text lives inside the chapter itself.  The Stacks Project uses one single shared preamble for
% every chapter; the snapshot stages it verbatim in sources/stacks-chow/preamble.tex.  The block below is
% copied mechanically, line for line, from that file's lines 14-18 (the Xy-pic setup) and lines 154-217
% (the native theorem declarations together with the mathematical macros).  Nothing was added, removed or
% redefined.  The wrapper class and the amsmath/amssymb/amsthm/graphicx base packages are supplied by the
% assembler, so the class/\externaldocument part of the shared preamble is not repeated here.
% For commutative diagrams we use Xy-pic
\usepackage[all]{xy}

% We use 2cell for 2-commutative diagrams.
\xyoption{2cell}
% Theorem environments.
%
\theoremstyle{plain}
\newtheorem{theorem}[subsection]{Theorem}
\newtheorem{proposition}[subsection]{Proposition}
\newtheorem{lemma}[subsection]{Lemma}

\theoremstyle{definition}
\newtheorem{definition}[subsection]{Definition}
\newtheorem{example}[subsection]{Example}
\newtheorem{exercise}[subsection]{Exercise}
\newtheorem{situation}[subsection]{Situation}

\theoremstyle{remark}
\newtheorem{remark}[subsection]{Remark}
\newtheorem{remarks}[subsection]{Remarks}

\numberwithin{equation}{subsection}

% Macros
%
\def\lim{\mathop{\mathrm{lim}}\nolimits}
\def\colim{\mathop{\mathrm{colim}}\nolimits}
\def\Spec{\mathop{\mathrm{Spec}}}
\def\Hom{\mathop{\mathrm{Hom}}\nolimits}
\def\Ext{\mathop{\mathrm{Ext}}\nolimits}
\def\SheafHom{\mathop{\mathcal{H}\!\mathit{om}}\nolimits}
\def\SheafExt{\mathop{\mathcal{E}\!\mathit{xt}}\nolimits}
\def\Sch{\mathit{Sch}}
\def\Mor{\mathop{\mathrm{Mor}}\nolimits}
\def\Ob{\mathop{\mathrm{Ob}}\nolimits}
\def\Sh{\mathop{\mathit{Sh}}\nolimits}
\def\NL{\mathop{N\!L}\nolimits}
\def\CH{\mathop{\mathrm{CH}}\nolimits}
\def\proetale{{pro\text{-}\acute{e}tale}}
\def\etale{{\acute{e}tale}}
\def\QCoh{\mathit{QCoh}}
\def\Ker{\mathop{\mathrm{Ker}}}
\def\Im{\mathop{\mathrm{Im}}}
\def\Coker{\mathop{\mathrm{Coker}}}
\def\Coim{\mathop{\mathrm{Coim}}}

% Boxtimes
%
\DeclareMathSymbol{\boxtimes}{\mathbin}{AMSa}{"02}

%
% Macros for moduli stacks/spaces
%
\def\QCohstack{\mathcal{QC}\!\mathit{oh}}
\def\Cohstack{\mathcal{C}\!\mathit{oh}}
\def\Spacesstack{\mathcal{S}\!\mathit{paces}}
\def\Quotfunctor{\mathrm{Quot}}
\def\Hilbfunctor{\mathrm{Hilb}}
\def\Curvesstack{\mathcal{C}\!\mathit{urves}}
\def\Polarizedstack{\mathcal{P}\!\mathit{olarized}}
\def\Complexesstack{\mathcal{C}\!\mathit{omplexes}}
% \Pic is the operator that assigns to X its picard group, usage \Pic(X)
% \Picardstack_{X/B} denotes the Picard stack of X over B
% \Picardfunctor_{X/B} denotes the Picard functor of X over B
\def\Pic{\mathop{\mathrm{Pic}}\nolimits}
\def\Picardstack{\mathcal{P}\!\mathit{ic}}
\def\Picardfunctor{\mathrm{Pic}}
\def\Deformationcategory{\mathcal{D}\!\mathit{ef}}

% --- appended 2026-10-12 (Chief): the project own definitions of the reference/slogan/history
% display environments, copied mechanically from the shared Stacks preamble lines 8-12
% (\usepackage{verbatim} plus the three \newenvironment lines, which the chapter preamble keeps
% verbatim from the project preamble). Nothing added, removed or redefined: these environments
% carry no mathematical content and the project itself typesets them as comments. Without these
% three lines a counted statement whose own text uses them fails with Environment slogan undefined.
\usepackage{verbatim}
\newenvironment{reference}{\comment}{\endcomment}
\newenvironment{slogan}{\comment}{\endcomment}
\newenvironment{history}{\comment}{\endcomment}
\usepackage{cleveref}
\crefname{originallemma}{original lemma}{original lemmas}
\Crefname{originallemma}{Original lemma}{Original lemmas}
\makeatletter
\newlabel{lemma-Noetherian-scheme-regular-ideal}{{063I}{source}{}{sourcepointer-1}{}}
\newlabel{lemma-Noetherian-scheme-regular-ideal@cref}{{[originallemma][0][]063I}{[1][0][]source}{}{}{}}
\makeatother
\crefname{originallemma}{original lemma}{original lemmas}
\Crefname{originallemma}{Original lemma}{Original lemmas}
\makeatletter
\newlabel{lemma-composition-regular-immersion}{{067Q}{source}{}{sourcepointer-2}{}}
\newlabel{lemma-composition-regular-immersion@cref}{{[originallemma][0][]067Q}{[1][0][]source}{}{}{}}
\makeatother
\crefname{originallemma}{original lemma}{original lemmas}
\Crefname{originallemma}{Original lemma}{Original lemmas}
\makeatletter
\newlabel{lemma-generate-regular-ideal}{{067N}{source}{}{sourcepointer-3}{}}
\newlabel{lemma-generate-regular-ideal@cref}{{[originallemma][0][]067N}{[1][0][]source}{}{}{}}
\makeatother
\crefname{originallemma}{original lemma}{original lemmas}
\Crefname{originallemma}{Original lemma}{Original lemmas}
\makeatletter
\newlabel{lemma-transitivity-conormal-quasi-regular}{{063N}{source}{}{sourcepointer-4}{}}
\newlabel{lemma-transitivity-conormal-quasi-regular@cref}{{[originallemma][0][]063N}{[1][0][]source}{}{}{}}
\makeatother
\title{Splitting of the conormal sequence for a composite of immersions detects regularity of both factors over a locally Noetherian base}
\author{The Stacks Project Authors}
\date{stacks-project:divisors.tex:local-snapshot-20261012; GNU FDL 1.2 or later}
\begin{document}
\maketitle
\noindent\textit{Editorial note.} All delivered mathematical text is the author's own, copied from the pinned Stacks Project chapter file pinned-original-source.tex; no statement or proof step was rewritten, reordered or removed.\par\medskip
\noindent\textit{Editorial note.} Following the multi-claim convention for this source, source\_url is the source-level URL https://stacks.math.columbia.edu/ and the source label lemma-extra-permanence-regular-immersion-noetherian is carried in claim\_key and here rather than in a per-claim tag URL.\par\medskip
\noindent\textit{Editorial note.} The preamble is the chapter's shared Stacks Project preamble, retained as sources/stacks-divisors/retained-preamble.tex: a line-for-line excerpt of the shared Stacks preamble staged with this snapshot. No macro was invented or redefined, and this card adds no preamble text of its own.\par\medskip
\noindent\textit{Editorial note.} Statement span 5144-5163 and proof span 5165-5197 come from the halfA pool card operations/pool-cards/stacks-divisors-halfA.json and were re-read in the pinned chapter file before this card was written.\par\medskip
\noindent\textit{Editorial note.} The statement contains a displayed three-term conormal sequence with a split exactness condition; it is copied verbatim from the counted statement span.\par\medskip
\noindent\textit{Editorial note.} The proof refers to 'all 4 types of regular immersions' agreeing in the locally Noetherian case; the author's wording is preserved.\par\medskip
\noindent\textit{Editorial note.} 4 cross-references occur inside the counted spans; they are listed in external\_references as navigation pointers only.\par\medskip
\noindent\textit{Allowed standard prerequisites (editorial).}\begin{itemize}
\item Over a locally Noetherian scheme all four notions of regular immersion agree and regularity can be checked on local rings, cited from the same chapter. Cited at proof lines 5165-5168. Reduces the whole statement to a statement about the local rings at the three points.
\item The composite of two regular immersions is a regular immersion, and the transitivity statement relating the conormal sequence of a composite to those of its factors, both cited from the same chapter. Cited at proof lines 5169-5172. Gives the implications from the first condition to the second and from the second to the third.
\item Any minimal sequence of generators of a quasi-regular ideal is a quasi-regular, hence regular, sequence, cited from the same chapter. Cited at proof lines 5180-5182. Makes a minimal generating set of the ideal of the composite immersion a regular sequence.
\item A minimal generating set of an ideal whose conormal sequence has a splitting compatible with a subideal can be chosen with the subideal's generators first, and a regular sequence splits into a regular sequence and a regular sequence modulo its first part. Used at proof lines 5183-5195. Produces the two shorter regular sequences that witness regularity of the individual immersions.
\end{itemize}\medskip
\ifdefined\hypertarget\hypertarget{sourcepointer-1}{}\fi
\noindent\textit{Source reference.} Original lemma 063I: focused-seeds/AGP/sources/stacks-snapshot-20261011/tags.txt, source lines 7832--7832. Same-chapter reference used in the retained author text at source line(s) 4939, 4970, 5169, 5234, 5705. It is supplied only as a navigation pointer so that the retained text compiles; the counted item delivers the author's own statement and the author's complete proof, and no part of the referenced item is copied in or substituted for the counted argument.\par\medskip
\ifdefined\hypertarget\hypertarget{sourcepointer-2}{}\fi
\noindent\textit{Source reference.} Original lemma 067Q: focused-seeds/AGP/sources/stacks-snapshot-20261011/tags.txt, source lines 7979--7979. Same-chapter reference used in the retained author text at source line(s) 5171, 5669, 5817. It is supplied only as a navigation pointer so that the retained text compiles; the counted item delivers the author's own statement and the author's complete proof, and no part of the referenced item is copied in or substituted for the counted argument.\par\medskip
\ifdefined\hypertarget\hypertarget{sourcepointer-3}{}\fi
\noindent\textit{Source reference.} Original lemma 067N: focused-seeds/AGP/sources/stacks-snapshot-20261011/tags.txt, source lines 7977--7977. Same-chapter reference used in the retained author text at source line(s) 5113, 5133, 5183. It is supplied only as a navigation pointer so that the retained text compiles; the counted item delivers the author's own statement and the author's complete proof, and no part of the referenced item is copied in or substituted for the counted argument.\par\medskip
\ifdefined\hypertarget\hypertarget{sourcepointer-4}{}\fi
\noindent\textit{Source reference.} Original lemma 063N: focused-seeds/AGP/sources/stacks-snapshot-20261011/tags.txt, source lines 7837--7837. Same-chapter reference used in the retained author text at source line(s) 5173. It is supplied only as a navigation pointer so that the retained text compiles; the counted item delivers the author's own statement and the author's complete proof, and no part of the referenced item is copied in or substituted for the counted argument.\par\medskip
\begin{lemma}
\label{lemma-extra-permanence-regular-immersion-noetherian}
Let $i : Z \to Y$ and $j : Y \to X$ be immersions of schemes.
Pick $z \in Z$ and denote $y \in Y$, $x \in X$ the corresponding points.
Assume $X$ is locally Noetherian.
The following are equivalent
\begin{enumerate}
\item $i$ is a regular immersion in a neighbourhood of $z$ and $j$
is a regular immersion in a neighbourhood of $y$,
\item $i$ and $j \circ i$ are regular immersions in a neighbourhood of $z$,
\item $j \circ i$ is a regular immersion in a neighbourhood of $z$ and the
conormal sequence
$$
0 \to i^*\mathcal{C}_{Y/X} \to
\mathcal{C}_{Z/X} \to
\mathcal{C}_{Z/Y} \to 0
$$
is split exact in a neighbourhood of $z$.
\end{enumerate}
\end{lemma}

\begin{proof}
Since $X$ (and hence $Y$) is locally Noetherian all 4 types of regular
immersions agree, and moreover we may check whether a morphism is a
regular immersion on the level of local rings, see
Lemma \ref{lemma-Noetherian-scheme-regular-ideal}.
The implication (1) $\Rightarrow$ (2) is
Lemma \ref{lemma-composition-regular-immersion}.
The implication (2) $\Rightarrow$ (3) is
Lemma \ref{lemma-transitivity-conormal-quasi-regular}.
Thus it suffices to prove that (3) implies (1).

\medskip\noindent
Assume (3). Set $A = \mathcal{O}_{X, x}$. Denote $I \subset A$ the kernel
of the surjective map $\mathcal{O}_{X, x} \to \mathcal{O}_{Y, y}$ and
denote $J \subset A$ the kernel
of the surjective map $\mathcal{O}_{X, x} \to \mathcal{O}_{Z, z}$.
Note that any minimal sequence of elements generating $J$ in $A$
is a quasi-regular hence regular sequence, see
Lemma \ref{lemma-generate-regular-ideal}.
By assumption the conormal sequence
$$
0 \to I/IJ \to J/J^2 \to J/(I + J^2) \to 0
$$
is split exact as a sequence of $A/J$-modules. Hence we can pick
a minimal system of generators $f_1, \ldots, f_n, g_1, \ldots, g_m$
of $J$ with $f_1, \ldots, f_n \in I$ a minimal system of generators of $I$.
As pointed out above $f_1, \ldots, f_n, g_1, \ldots, g_m$ is a regular
sequence in $A$. It follows directly from the definition of a regular
sequence that $f_1, \ldots, f_n$ is a regular sequence in $A$ and
$\overline{g}_1, \ldots, \overline{g}_m$ is a regular sequence in
$A/I$. Thus $j$ is a regular immersion at $y$ and $i$ is a regular
immersion at $z$.
\end{proof}

\end{document}
